\def\stat{bystrov}





\def\tit{ОСНОВЫ ЖИВУЧЕСТИ АВТОМАТИЗИРОВАННЫХ ОРГАНИЗАЦИЙ  


ИНФОРМАЦИОННОГО ОБЩЕСТВА}





\def\titkol{Основы живучести автоматизированных организаций  


информационного общества}








\def\aut{И.\,И.~Быстров$^1$, В.\,Н.~Веселов$^2$, К.\,К.~Колин$^3$}





\def\autkol{И.\,И.~Быстров, В.\,Н.~Веселов, К.\,К.~Колин}





\titel{\tit}{\aut}{\autkol}{\titkol}





\index{Быстров И.\,И.}


\index{Веселов В.\,Н.}


\index{Колин К.\,К.}


\index{Bystrov I.\,I.}


\index{Veselov V.\,N.}


\index{Kolin K.\,K.}








%{\renewcommand{\thefootnote}{\fnsymbol{footnote}}


%\footnotetext[1]


%{Работа выполнена при поддержке РФФИ (проект 15-07-02053).}}





\renewcommand{\thefootnote}{\arabic{footnote}}


\footnotetext[1]{Институт проблем информатики Федерального исследовательского центра 


<<Информатика и~управление>> 


Российской академии наук, \mbox{ibystrov@ipiran.ru}}


\footnotetext[2]{Институт проблем информатики Федерального исследовательского центра 


<<Информатика и~управление>>


Российской академии наук, \mbox{vveselov@ipiran.ru}}


\footnotetext[3]{Институт проблем информатики Федерального исследовательского центра 


<<Информатика и~управление>>


Российской академии наук, \mbox{kolinkk@mail.ru}}








 \label{st\stat}





                  \thispagestyle{headings}


                  


  \vspace*{-14pt}


  





  \Abst{Рассматриваются вопросы, связанные с~обеспечением живучести информационного 


общества, в~котором широко используются ин\-фор\-ма\-ци\-он\-но-те\-ле\-ком\-му\-ни\-ка\-ци\-он\-ные 


сис\-те\-мы и~технологии. Ставится проблема создания методологии проектирования 


комплексной сис\-те\-мы обеспечения живучести (КСОЖ) 


автоматизированных организаций. 


Рассматриваются модели и~показатели обеспечения живучести 


ин\-фор\-ма\-ци\-он\-но-те\-ле\-ком\-му\-ни\-ка\-ци\-он\-ных сис\-тем.}


  


  \KW{информационное общество; информационно-те\-ле\-ком\-му\-ни\-ка\-ци\-он\-ная система; 


живучесть автоматизированных организаций; модели и~показатели обеспечения жи\-ву\-чести}





\DOI{10.14357\!/\!08696527180411} 








\vskip 10pt plus 9pt minus 6pt








      \thispagestyle{headings}


      


     % \vspace*{-10pt} 


     


     \vspace*{-20pt}


  


  \section{Актуальность проблемы}


  


  На сегодняшний день проблемы, связанные с~информатизацией общества, 


приобрели новую актуальную значимость, так как уровень применения сетевых 


информационных технологий (ИТ)  


и~ин\-фор\-ма\-ци\-он\-но-ком\-му\-ни\-ка\-ци\-он\-ных технологий (ИКТ), 


который является одним из важных факторов конкурентности страны 


в~мировой экономике, а также в~политической и~военной сфере, в~настоящее 


время становится ма\-те\-ри\-аль\-но-тех\-ни\-че\-ской базой цифровой 


экономики.


  


  Комплексная автоматизация общества предполагает сис\-тем\-ный подход, 


широкое внедрение средств вычислительной техники, коммуникаций 


и~соответствующего программного обеспечения в~деятельность 


автоматизированных организаций, т.\,е.\ организаций, деятельность которых 


осуществляется на основе автоматизации функциональных процессов. 


В~результате этого значительная часть функциональных процессов, 


определяющих целевое предназначение этих организаций, выполняется 


с~использованием ин\-фор\-ма\-ци\-он\-но-те\-ле\-ком\-му\-ни\-ка\-ци\-он\-ных 


сис\-тем (ИТС).


  


  При этом если ранее в~случае выхода из строя автономных компьютерных 


сис\-тем организации сравнительно легко переходили на ручные методы 


управления, то при комплексной автоматизации такой подход стал практически 


невозможен. Это привело к~тому, что жизнедеятельность автоматизированной 


организации стала зависеть от состояния ИТС.


  


  В последние годы все большее развитие и~применение, особенно в~крупных 


государственных, коммерческих и~военных организациях, стали получать 


современные распределенные по всей структуре общества и~государства  


ин\-фор\-ма\-ци\-он\-но-те\-ле\-ком\-му\-ни\-ка\-ци\-он\-ные сис\-те\-мы (сети), 


представляющие собой интеграцию информационных и~коммуникационных 


составляющих различных компьютерных сис\-тем (сетей) с~опорой на 


взаимосвязанную совокупность крупных многофункциональных центров 


коллективной обработки информации.


  


  Внедрение этих сис\-тем существенно повысило эффективность применения 


ИТ в~деятельности общества, в~формировании и~развитии национальной 


цифровой экономики, обеспечении интересов личности и~реализации 


стратегических национальных приоритетов~[1]. В~связи с~этим приобретает 


особую важность и~актуальность проблема \textit{обеспечения устойчивого 


и~бесперебойного функционирования автоматизированных субъектов 


и~объектов общества и~государства на базе ИТС}, их способности надежно 


сохранять и~восстанавливать свою функциональность в~условиях воздействия 


всего спектра угроз.


  


  Современные ИТС представляют собой \textit{организационно-технические 


структуры} с~высокой концентрацией технических, программных средств 


и~информации, которые могут быть объектами случайных или 


преднамеренных злоумышленных воздействий на основе использования 


свойства уязвимости информации, инфраструктуры сис\-тем (сетей), под\-сис\-тем, 


процессов и~оборудования.


  


  Структурная сложность и~территориальная распределенность ИТС и~их 


компонентов обусловливают их уязвимость от множества внутренних 


и~внешних угроз различной природы, что не позволяет превентивными мерами 


обеспечить их гарантированную защиту. Нарушение их функциональности 


приводит к~возникновению высоких рисков в~жизнедеятельности субъектов 


и~объектов общества, к~нарушению их работоспособности и~информационной 


безопасности, а~в~некоторых случаях может приводить и~к~остановке 


деятельности критических объектов инфраструктуры общества и~государства.


  


  Наличие свойства уязвимости ИТС позволяет реализовать угрозы нарушения 


конфиденциальности и~целостности информации, обойти превентивные меры 


защиты компонентов сис\-тем и~сетей, нарушить их функциональные процессы 


и~физическое состояние. Так, в~2015~г.\ в~России было 


выявлено~189~уязвимостей в~компонентах автоматизированных систем управ\-ле\-ния, большинство из которых 


являлись критическими (49\%) или имели средний уровень опасности (42\%).


  


  Все это приводит к~тому, что определенные субъекты (коалиции, 


государства, организации и~личности) стремятся через инфраструктуру ИТС, 


компьютерные сис\-те\-мы (сети, Интернет) осуществлять \textit{вредоносные 


кибервоздействия} на жизнедеятельность общества и~государства в~своих 


интересах.


  


  Ввиду невозможности обеспечить достаточно высокую киберзащиту 


общества одной из важнейших задач обеспечения национальной безопасности 


страны становится придание автоматизированным организациям и~всей 


инфраструктуре государства способности непрерывно функционировать 


в~условиях воздействия существующих угроз и~выхода из строя отдельных 


критичных объ\-ектов.


{\looseness=1





}


  


  Согласно результатам одного из последних исследований, проведенных 


компанией Forester, планирование и~управление непрерывностью бизнеса, 


а~также обеспечение восстановления деятельности в~непредвиденных 


обстоятельствах и~при стихийных бедствиях определяются как два 


приоритетных направления деятельности организаций, которые сегодня 


находятся в~центре внимания руководителей ведущих финансовых и~других 


автоматизированных организаций на всех континентах~\cite{2-bys}.


  


  При таком подходе речь должна идти о новом комплексном свойстве~--- 


\textit{живучести} (ГОСТ~34.003-90) как отдельных автоматизированных 


организаций, так и~живучести всего информационного общества в~целом.


  


  \section{Живучесть информационного общества: содержание проблемы 


и~методология ее решения}


  


  Под \textit{живучестью информационного общества} будем понимать такое 


его состояние, при котором оно способно в~условиях воздействия угроз 


поддерживать свою жизнедеятельность, доступность услуг, а также сохранять 


и~восстанавливать функциональность критичных автоматизированных 


организаций с~минимизацией ущерба информационной и~национальной 


безопасности.


  


  Проблему живучести информационного общества как в~теоретическом, так 


и~в~практическом плане целесообразно решать комплексно. Это должно 


осуществляться на базе \textit{об\-ще\-сис\-тем\-ной методологии}, объединяющей 


на сис\-тем\-ном уровне правовые основы Стратегии развития информационного 


общества в~Российской Федерации на~2017--2030~гг.\ и~Доктрины 


информационной безопасности РФ~\cite{3-bys}, а также последние научные 


достижения множества дисциплин, связанных с~применением ИТ и~их 


реализацией в~виде ИТС.


  


  Ключевым моментом методологии обеспечения живучести 


информатизированного общества является создание \textit{теории и~практики 


построения единого технологического и~организационного поля}, в~котором 


критичные сис\-те\-мы, ИТС и~автоматизированные организации общества 


и~государства сохраняют свою способность противодействовать угрозам, 


осуществлять предоставление ИТ-услуг и~восстанавливать свою 


функциональную деятельность.


  


  В настоящее время сложилось несколько направлений решения проблемы 


обеспечения информационной безопасности общества с~использованием 


средств информатизации. Представители первого направления связывают 


решение этой проблемы с~акцентом на \textit{техническую защиту 


информации} в~сис\-те\-мах, без их адаптации к~угрозам технического 


и~физического разрушения инфраструктуры, возможным рискам 


и~последствиям их для общества и~государства.


  


  Другие рассматривают проблему влияния устойчивости сис\-тем на 


жиз\-не\-де\-ятель\-ность общества только с~точки зрения \textit{надежности 


технических и~про\-граммных средств}.


  


  Третьи предлагают решать эту проблему с~позиции противодействия 


угрозам, связанным с~манипулированием функциональностью сис\-те\-мы, 


ресурсами, информацией, блокированием и~физическим разрушением 


инфраструктуры и~киберсреды объектов общества, органов военного 


и~государственного управления.


  


  По каждому из этих направлений в~настоящее время выработаны 


определенные требования к~безопасности общества и~государства, но это 


в~основном рекомендации <<частного>>, корпоративного, технического 


и~правового плана. Несмотря на то что эти подходы тесно связаны между 


собой, их раздельное применение решает проблемы информационной 


безопасности общества и~национальной безопасности не в~полной мере.


  


  Живучесть, по мнению авторов, является \textit{комплексным свойством} 


автоматизированного общества и~государства и~зависит от множества 


факторов. Среди множества свойств, которые целесообразно учитывать, 


исследуя жизнеспособность информатизированного общества и~государства, 


необходимо выделять их способность существовать и~сохранять свою 


функциональность на основе \textit{живучести критичных 


автоматизированных организаций, функционирующих} на базе  


ИТС~\cite{4-bys}.


  


  Понятие <<\textit{живучесть автоматизированной организации>>} 


в~традиционном понимании~--- это функциональное свойство устойчивости 


сложных сис\-тем, каковыми являются современные автоматизированные 


субъекты информационного общества и~государства.


  


  В качестве нового определения под \textit{живучестью 


автоматизированной организации будем понимать такое ее свойство, при 


котором компьютерные сис\-те\-мы (сети), интегрированные в~инфраструктуру 


ИТС, способны поддерживать бесперебойное функционирование деятельности 


организации, доступность ИТ-услуг, сохраняя и~восстанавливая свою 


функциональность с~использованием специальных механизмов адаптации 


к~условиям воздействия угроз}.


  


  Все угрозы для функционирования ИТС могут быть классифицированы по 


целому ряду критериев. Базельский комитет предлагает следующую 


классификацию источников этих угроз:


  \begin{itemize}


  \item несанкционированные действия персонала организации и~некорректное 


выполнение своих обязанностей;


  \item нарушения в~результате сбоя функционирования автоматизированных 


сис\-тем и~процессов доступа к~ИТ-сервисам;


  \item преднамеренные угрозы и~действия, исходящие из внешней среды.


  \end{itemize}


  


  Наряду с~традиционными угрозами, возникающими в~инфраструктуре ИТС, 


в~последние годы с~распространением глобальных, национальных 


корпоративных ИТС крупных организаций определилась техническая среда, 


через которую осуществляются кибервойны, кибератаки криминальных 


элементов общества, террористов, информационно-технологическое 


воздействие на личности и~принимаемые ими решения~\cite{5-bys}. 


В~результате этого в~зависимости от сферы применения ИТС, их уязвимости 


и~характера угроз создаются опасности для жизнедеятельности организаций 


и~всего общества в~целом, которые по своей общей на\-прав\-лен\-ности 


в~соответствии с~Доктриной информационной безопасности оказывают 


влияние:


  \begin{itemize}


  \item на конституционные права и~свободы граждан, реализуемые 


в~информационной сфере;


  \item состояние инфраструктуры;


  \item информационные ресурсы.


  \end{itemize}


  


  История проблемы живучести деятельности автоматизированных 


организаций началась с~момента появления электронных


вычислительных машин (ЭВМ). По мере расширения сферы 


использования ЭВМ с~появлением автоматизированных информационных 


сис\-тем и~центров обработки данных (ЦОД) возникла проблема защиты 


информации в~сис\-те\-мах, расширился круг решаемых задач, а~деятельность 


организаций стала все больше зависеть от качества и~непрерывности 


функционирования этих сис\-тем и~центров. В~1980-х~гг.\ окончательно 


сформировалась проблема аварийного восстановления в~случае возникновения 


угроз, способных нарушить функционирование сис\-тем, и~в~первую очередь 


ЦОД.


  


  В 1990-е гг.\ на смену автономным ЦОД пришли персональные 


компьютеры и, соответственно, изменился подход к~резервированию данных 


и~аварийному восстановлению автоматизированных информационных сис\-тем. 


Большинство организаций отказались от централизованных мейнфреймов, 


заменяя их большим числом серверов и~пользовательских автоматизированных 


мест, распределенных по всей структуре организации. Возникло множество 


достаточно автономных, разнородных автоматизированных сис\-тем 


документальной связи, различных информационных и~ин\-фор\-ма\-ци\-он\-но-управ\-ля\-ющих сис\-тем.


  


  На смену экстенсивному развитию в~конце 1990-х гг.\ пришло интенсивное 


развитие в~направлении создания крупномасштабных мультисервисных ИТС, 


интегрирующих информационные и~телекоммуникационные ресурсы 


различных отдельных сис\-тем и~сетей. Эти сис\-те\-мы органически слились со 


структурой организаций и~стали определять значительную долю процессов их 


де\-ятель\-ности, скорость их выполнения, и~непрерывность процессов стала 


зависеть от бесперебойной обработки информации, ее хранения и~качества 


сервисов и~ИТ-услуг, предоставляемых ИТС. В~связи с~этим на смену термину 


<<аварийное восстановление>> пришел термин <<непрерывность бизнеса>>.


  


  В последние 10--15~лет ведущие мировые компании, в~первую очередь 


финансового сектора, в~условиях постоянно растущих угроз информационным 


и~коммуникационным технологиям и~частоты возникновения чрезвычайных 


ситуаций осознали степень зависимости бизнеса от состояния непрерывности 


функционирования и~доступности сервисов ИТС. Поэтому начался процесс 


целенаправленного внедрения требований непрерывности 


деятельности~организаций на базе ИТС.


  


  В США, Великобритании, Австралии, Новой Зеландии, Японии, Сингапуре 


в~эти годы создаются институты, которые разрабатывают стандарты 


и~рекомендации в~области <<непрерывности функционирования и~доступности 


сервисов ИТС>>. В~результате этого понятие \textit{непрерывности 


функционирования и~доступности ИТ-услуг} стало базовым при решении 


проблем живучести автоматизированных организаций.


  


  \section{Методология обеспечения живучести автоматизированных 


организаций}


  


  Главная цель обеспечения живучести автоматизированной организации, 


функционирующей на базе ИТС, заключается в~сохранении ее 


работоспособности в~случае полного или частичного нарушения  


ре\-сурс\-но-сер\-вис\-ных возможностей ИТС.


 % 


  Здесь важно учитывать, что в~случае нарушения этих требований ИТС 


превращается из незаменимого средства технической поддержки штатной 


работы в~средство, которое может парализовать всю деятельность организации. 


Отсюда возникает центральное требование к~ИТС~--- \textit{способности 


адаптироваться к~вероятным угрозам и~сохранять свою способность 


к~функционированию}. Для достижения этих качеств ИТС должна обладать 


комплексным свойством живучести, которое предполагает ее способности:


  \begin{itemize}


  \item обеспечения непрерывности функционирования и~доступности 


сервисов и~ИТ-услуг;


  \item обеспечения комплексной безопасности инфраструктуры 


и~информационных ресурсов;


  \item управление рисками, порождаемыми самой ИТС.


  \end{itemize}


  


  Основой создания комплексной сис\-те\-мы обеспечения технической 


жи\-ву\-чести автоматизированной организации должна стать 


\textit{об\-ще\-сис\-тем\-ная методология}, интегрирующая в~себе  


ор\-га\-ни\-за\-ци\-он\-но-тех\-ни\-че\-ский инструментарий реализации целей 


и~задач автоматизированной организации. Об\-ще\-сис\-тем\-ная методология 


обеспечения живучести автоматизированной организации на базе ИТС 


охватывает широкий круг вопросов, в~частности:


  \begin{itemize}


  \item разработку научных основ обеспечения устойчивости, применение 


методов сис\-тем\-но\-го и~процессного подхода;


  \item создание моделей формирования и~поддержки ИТ-услуг в~различных 


условиях обстановки с~ориентацией на пользователей;


  \item разработку принципов проектирования, комплексирования 


и~интеграции методов и~средств в~инфраструктуру ИТС и~ее технологических 


процессов;


  \item создание методов и~способов целостного управления с~основных 


и~резервных площадок с~необходимым уровнем информационной поддержки 


процессов принятия решений в~чрезвычайных ситуациях;


  \item создание методики определения критически важных объектов, сис\-тем, 


сервисов и~данных, оценки технических рисков и~приоритетов восстановления 


компонентов ИТС, подвергшихся разрушению;


  \item разработку принципов построения ре\-сурс\-но-сер\-вис\-ной сис\-те\-мы 


эксплуатации для полной поддержки живучести действующих ИТС 


и~обеспечения предоставления пользователям критически важных сервисов 


и~ИТ-услуг в~период действия чрезвычайной ситуации;


  \item разработку нормативно-методической базы, планов, регулирующих 


вопросы обеспечения живучести организации на базе ИТС и~восстановления 


сервисов в~случае возникновения угроз и~чрезвычайных ситуаций.


  \end{itemize}


  


  При решении этих вопросов должны проводиться:


  \begin{itemize}


  \item анализ угроз, приводящих к~нарушению живучести ИТС;


  \item оценка уязвимости ИТС, ее компонентов и~процессов;


  \item анализ рисков для ИТС, порождаемых угрозами, и~оценка их влияния на 


живучесть, надежность и~помехозащищенность ап\-па\-рат\-но-про\-грам\-мных 


средств;


  \item определение критичности объектов, сис\-тем (сетей) и~под\-сис\-тем, 


нарушение функционирования которых оказывает влияние на способность ИТС 


выполнять свои функции в~заданных режимах;


  \item анализ и~оценка ущерба организации, возникающего в~случае 


нарушения непрерывности функционирования и~доступности сервисов (НФДС) 


ИТС.


  \end{itemize}


  


  Исходя из приведенного выше определения живучести, можно дать 


следующее определение общего требования к~живучести функционирования 


любого автоматизированного субъекта общества и~государства на базе ИТС.


  


  {\bfseries\textit{Живучесть автоматизированного субъекта общества 


и~государства}} представляет собой стратегическую, оперативную 


и~тактическую способность планировать и~осуществлять свою работу в~случае 


возникновения негативных факторов, нарушающих его деятельность, 


направленную на обеспечение непрерывности и~комплексной безопасности 


деловых операций в~соответствии с~установленным приемлемым уровнем 


рисков для ИТС.


  


  Деятельность субъекта общества (организации) на каждом уровне следует 


рассматривать как комплексный сбалансированный автоматизированный 


процесс, охватывающий такие институциональные аспекты, как правовая база, 


стандарты, рыночная конъюнктура, механизмы стратегического планирования 


процессов обеспечения непрерывности и~восстановления деятельности 


(ОНиВД), а также информационные и~технические аспекты, включая 


требования к~способности ИТС обеспечивать функциональные процессы 


организации.


  


  \textit{На стратегическом уровне} определяются принципы деятельности 


организации на базе ИТС, живучесть автоматизированной организации 


и~комплекс ор\-га\-ни\-за\-ци\-он\-но-тех\-ни\-че\-ских мер реализации 


стратегии и~плана действий в~чрезвычайных ситуациях по защите объектов, 


персонала, критических функциональных процессов.


  


  \textit{На оперативном уровне} определяются концептуальные основы создания КСОЖ:


  \begin{itemize}


  \item идентификация критически важных для организации процессов 


и~поддерживающих ресурсов ИТС;


  \item оценка рисков и~возможного ущерба от нарушения функционирования 


ИТС;


  \item определение требований к~обеспечению НФДС ИТС;


  \item выбор общих способов применения превентивных мер обеспечения 


безопасности;


  \item выбор методов устранения всех видов рисков, которые не удалось 


исключить с~помощью превентивных мер, а также мер, предоставляемых на 


период восстановления;


  \item разработка требований, синхронизирующих подходы к~управлению 


непрерывностью функционирования организации и~ИТС;


  \item разработки методов и~средств реализации требований нормативных 


документов и~методик оценки соответствия организационных и~технических 


решений создания КСОЖ.


  \end{itemize}


  


  \textit{На тактическом уровне} определяется организационно-техническое 


обеспечение и~поддержка процессов практического проектирования КСОЖ 


ИТС, решаются проблемы технических решений НФДС, эксплуатационной 


живучести, безопасности ресурсов и~информации, а также восстановления 


функционирования ИТС в~приемлемый для организации период времени.


  


  Основная цель создания КСОЖ ИТС~--- техническая поддержка нормальной 


деятельности организации и~восстановление нарушенной работоспособности 


ИТС в~максимально короткие сроки с~минимизацией оперативного влияния на 


ее процессы. Реализация мероприятий в~организации на каждом уровне 


определяется стратегией живучести общества и~государства, учитывающей 


необходимый баланс между экономическими, финансовыми, деловыми, 


эксплуатационными рисками и~планированием мероприятий на случай 


чрезвычайных происшествий.


{\looseness=1





}


  


  Организационно-технической базой восстановления функционирования ИТС 


являются критерии:


  \begin{itemize}


  \item быстрого и~своевременного восстановления функциональности;


  \item высокого уровня уверенности в~том, что (за счет повседневного 


использования и~проверочного тестирования) внутренние и~внешние 


мероприятия по обеспечению непрерывности процессов эффективны 


и~совместимы.


  \end{itemize}


  


  При этом КСОЖ должна удовлетворять:


  \begin{description}


  \item[\,] \textit{требованиям к~времени восстановления}:


  \begin{itemize}


  \item  для ведущих критичных организаций~--- процессы восстановления 


и~возобновления своей деятельности в~течение дня с~момента аварии;


  \item для организаций, играющих важную роль для других критически 


важных организаций,~--- возобновление деятельности в~пределах нескольких 


часов;


  \end{itemize}


  


 \item[\,] \textit{требованиям к~структурной устойчивости}:


  \begin{itemize}


  \item иметь для восстановления после широкомасштабной катастрофы 


в~заданное время географически разнесенные основную и~резервные площадки 


для выполнения деловых операций и~работы ЦОД;


  \item резервные площадки ИТС не должны зависеть от тех же 


инфраструктурных компонентов, что и~основная, включая средства 


транспортировки, телекоммуникации, водоснабжения, электроснабжения и~т.\,п.;


  \end{itemize}


  


\item[\,]  \textit{требованию к~готовности}~--- наличию эффективных планов 


обеспечения непрерывности и~восстановления деятельности организации 


в~условиях возникновения катастрофы.


\end{description}


  


  \section{Заключение}


  


  Развитие ИТС определяет процесс создания на\-уч\-но-тех\-ни\-че\-ских 


и~экономических основ нового типа субъектов цифрового общества 


и~государства, у~которых вся их деятельность в~существенной степени зависит 


от состояния ИТС.


  


  Существующие перспективные ИТС и~технологии аккумулируют 


и~обрабатывают большие объемы критической информации и~могут 


практически контролировать и~регулировать деятельность любого субъекта 


общества и~государства. Это приводит к~тому, что выживание любого субъекта, 


общества и~государства становится зависимым от состояния ИТС и~ИТ, а~сами 


сис\-те\-мы стали объектом воздействия различного рода угроз и~влиять на 


жизнеспособность общества и~государства.


  


  В условиях существования информатизированного общества и~государства 


в~среде современных реальных и~потенциальных угроз ИТС, когда методы 


и~средства защиты будут отставать от процессов появления новых угроз, 


а~последствия от нарушения функциональности ИТС будут постоянно 


возрастать, \textit{проблема обеспечения их живучести} приобретает особую 


важность и~актуальность для национальной безопасности Российской 


Федерации.


  


  В настоящее время эта проблема еще не нашла должного научного 


исследования как многодисциплинарная комплексная проблема. В~составе 


этой проблемы необходимо в~первую очередь выделять ее технологические 


аспекты:


  \begin{itemize}


  \item вопросы категорирования критически важных объектов, сис\-тем 


и~сервисов;


  \item вопросы обеспечения комплексной безопасности ИТС, резервирования 


компонентов сис\-те\-мы, требований к~восстановлению и~эксплуатационной 


поддержке живучести, к~устойчивости функциональных процессов;


  \item вопросы создания комплексной сис\-те\-мы обеспечения живучести ИТС.


  \end{itemize}


  


  Однако в~дальнейшем необходимо будет развернуть комплексные 


междисциплинарные исследования этой проблемы, которые будут включать 


и~другие, в~том числе гуманитарные, аспекты этой проблемы. В~их числе:


  \begin{itemize}


  \item вопросы управления непрерывностью и~операционными рисками, 


по\-рож\-да\-емы\-ми ИТС, для жизнедеятельности общества и~государства;


  \item вопросы защиты материальных и~финансовых ресурсов, репутации на 


рынке, защиты сотрудников, инфраструктуры субъектов общества 


и~государства от нарушений и~рисков, порождаемых киберугрозами.


  \end{itemize}


  


  В~условиях развития цифровой экономики и~становления глобального 


информационного общества зависимость безопасности жизнедеятельности 


общества от качества функционирования ИТС и~уровня живучести 


автоматизированных организаций будет, безусловно, возрастать. Поэтому 


научные исследования и~прикладные разработки, направленные на решение 


этой актуальной проблемы, должны стать важным направлением мероприятий 


по укреплению национальной безопасности нашей страны~\cite{6-bys}.


  


  


 {\small\frenchspacing


{%\baselineskip=10.8pt


\addcontentsline{toc}{section}{Литература}


\begin{thebibliography}{9}


  \bibitem{1-bys}


  \Au{Колин К.\,К.} Информатизация общества и~глобализация.~--- Красноярск: СФУ, 2011. 


52~с.


  \bibitem{2-bys}


  Доклад компании Forrester Report Sun Gard от 02.09.2010.


  \bibitem{3-bys}


  Доктрина информационной безопасности Российской Федерации: Утверждена Указом 


Президента РФ от 5~декабря 2016~г. №\,646.


  \bibitem{4-bys}


  \Au{Быстров И.\,И.} Живучесть автоматизированных организаций.~--- М.: Майор, 2016. 


506~с.


  \bibitem{5-bys}


  \Au{Колин К.\,К.} Технологическое общество: глобальные тенденции, вызовы 


  и~угрозы~/\!\!/ Стратегические приоритеты, 2017. №\,1. С.~4--15.


  \bibitem{6-bys}


  \Au{Соколов И.\,А., Колин К.\,К.} Развитие информационного общества в~России 


и~актуальные проблемы информационной безопасности~/\!\!/ Информационное общество, 


2009. №\,4-5. С.~98--107.


        \end{thebibliography}


} }








\vspace*{-3pt}





\hfill{\small\textit{Поступила в~редакцию 05.06.18}}














%\vspace*{8pt}





%\hrule





%\vspace*{2pt}





\newpage





\vspace*{-28pt}





%\hrule





%\vspace*{9pt}











\def\tit{INFORMATION SOCIETY AUTOMATED ORGANIZATIONS  


SURVIVABILITY BASICS}





\def\titkol{Information society automated organizations  


survivability basics}








\def\aut{I.\,I.~Bystrov, V.\,N.~Veselov, and~K.\,K.~Kolin}


\def\autkol{I.\,I.~Bystrov, V.\,N.~Veselov, and~K.\,K.~Kolin}











\titel{\tit}{\aut}{\autkol}{\titkol}





\vspace*{-9pt}





\noindent


      Institute of Informatics Problems, Federal Research Center ``Computer 


Science and~Control'' of the Russian Academy of Sciences, 44-2~Vavilov Str., 


Moscow 119333, Russian Federation








 \noindent





\def\leftfootline{\small{\textbf{\thepage}


\hfill Sistemy i Sredstva Informatiki~---


Systems and Means of Informatics 2018 vol~28 no\,4}


}%


 \def\rightfootline{\small{Sistemy i Sredstva Informatiki~---


 Systems and Means of Informatics   2018   vol~28   no\,4


\hfill \textbf{\thepage}}}


  


    


  \Abste{The paper deals with the issues related to ensuring survivability of the information 


society in which information and telecommunication systems and technologies are widely used. The 


problem of creating a methodology for designing an integrated system for ensuring survivability of 


automated organizations is raised. Models and indicators of ensuring survivability of information 


and telecommunication systems are considered.}


  


  \KWE{information society; information-telecommunication system; automated organizations 


survivability; models and indicators of ensuring survivability}


  


\DOI{10.14357\!/\!08696527180411} 





%\vspace*{-24pt}





%\Ack


%\noindent





 








\renewcommand{\bibname}{\protect\rmfamily References}


%\renewcommand{\bibname}{\large\protect\rm References}





{\small\frenchspacing


{%\baselineskip=10.8pt


\addcontentsline{toc}{section}{References}


\begin{thebibliography}{9}


  \bibitem{1-bys-1}


  \Aue{Kolin, K.\,K.} 2011. \textit{Informatizatsiya obshchestva i~globalizatsiya} [Society 


informatization and globalization] . Krasnoyarsk: SFU Publs. 52~p.


  \bibitem{2-bys-1}


  Doklad kompanii Forrester [Company Forrester Report]. 02.09.2010. Sun Gard.


  \bibitem{3-bys-1}


  Doktrina informatsionnoy bezopasnosti Rossiyskoy Federatsii [Doctrine of information security 


of the Russian Federation]. 2016. Utverzhdena Ukazom Prezidenta RF ot 5~dekabrya 2016~g. 


No.\,646 [Approved by Decree of the President of RF of December~5, 2016, No.\,646]. 


  \bibitem{4-bys-1}


  \Aue{Bystrov, I.\,I.} 2016. \textit{Zhivuchest' avtomatizirovannykh organizatsiy} [Automated 


organization survivability]. Moscow: Mayor Publs. 506~p.


  \bibitem{5-bys-1}


  \Aue{Kolin, K.\,K.} 2017. Tekhnologicheskoe obshchestvo: global'nye tendentsii, vyzovy 


i~ugrozy [Technological society: Global tendentions, calls, and threats]. 


\textit{Strategic Priorities} 13(1):4--15.


  \bibitem{6-bys-1}


  \Aue{Sokolov, I.\,A., and K.\,K.~Kolin.} 2009. Razvitie informatsionnogo obshchestva v~Rossii 


i~aktual'nye problemy informatsionnoy bezopasnosti [Information society development in Russia 


and information security actual problems]. \textit{Informatsionnoe obshchestvo} [Information 


Society] 21(4-5):98--107.


\end{thebibliography}


} }








\vspace*{-3pt}





\hfill{\small\textit{Received June 5, 2018}}  


  


  \Contr


  


  


  \noindent


  \textbf{Bystrov Igor I.} (b.\ 1931)~--- Doctor of Science in technology, professor, 


  principal 


scientist, Institute of Informatics Problem, Federal Research Center ``Computer Science and 


Control'' of the Russian Academy of Sciences, 44-2~Vavilov Str., Moscow 119333; 


\mbox{ibystrov@ipiran.ru}


  


  \vspace*{5pt}


  


  \noindent


  \textbf{Veselov Vitaliy N.} (b.\ 1940)~--- Candidate of Science (PhD) in technology, leading 


engineer, Institute of Informatics Problem of the Federal Research Center ``Computer Science and 


Control'' of the Russian Academy of Sciences, 44-2~Vavilov Str., Moscow 119333,


Russian Federation; 


\mbox{vveselov@ipiran.ru}


  


  \vspace*{5pt}





  \noindent


  \textbf{Kolin Konstantin K.} (b.\ 1935)~--- Doctor of Science in technology, professor, 


Honored Scientist of the Russia Federation, principal scientist, Institute of Informatics Problem, 


Federal Research Center ``Computer Science and Control'' of the Russian Academy of Sciences, 


44-2~Vavilov Str., Moscow 119333, Russian Federation; \mbox{kolinkk@mail.ru}











\label{end\stat}








\renewcommand{\bibname}{\protect\rm Литература}


